\documentclass{article}
\usepackage{fontspec}
\setmainfont{DejaVu Serif}
\setsansfont{DejaVu Sans}
\setmonofont{DejaVu Sans Mono}
\usepackage{amsmath, amssymb}
\usepackage{geometry}
\usepackage{float}
\usepackage{url}
\geometry{a4paper, margin=1in}
\setlength{\emergencystretch}{3em}

\newcommand{\Overlap}{\operatorname{Overlap}}
\newcommand{\Atomic}{\operatorname{AtomicSlots}}
\newcommand{\Conf}{\operatorname{Conf}}
\newcommand{\Cost}{\operatorname{Cost}}

\begin{document}

\title{\textbf{A SAT Encoding for Fixed-Workbook Personal Timetable Selection}}
\author{Formal Method Notes}
\date{}
\maketitle

\section{Notation}

\subsection*{Input Objects}
\begin{itemize}
    \item $\mathcal{S}$: the set of sessions extracted from the workbook. Each session corresponds to one workbook row.
    \item $\mathcal{P} = \{1, \dots, m\}$: the set of user-selected teaching-team
    requests. Each request is $(course\_code, teaching\_team\_key)$.
    \item $\mathcal{C}_p$: the set of candidate LHPs for request $p \in \mathcal{P}$.
    \item $\Sigma_{p,c} \subseteq \mathcal{S}$: the set of sessions that belong to candidate $c \in \mathcal{C}_p$.
    \item $t(s)$: the fixed workbook timeslot of session $s$.
    \item $r(s)$: the room associated with session $s$.
    \item $k_{\max}$: the maximum number of solutions returned by the desktop command. The current default is $5$.
\end{itemize}

\subsection*{Time Domain}
Each timeslot has the form $(day, period)$. Period values are interpreted through atomic teaching blocks:
\begin{align*}
\Atomic("1") &= \{1\}, & \Atomic("2") &= \{2\}, \\
\Atomic("3") &= \{3\}, & \Atomic("4") &= \{4\}, \\
\Atomic("1-2") &= \{1,2\}, & \Atomic("2-3") &= \{2,3\}, \\
\Atomic("3-4") &= \{3,4\}, \\
\Atomic(\text{Sáng}) &= \{1,2\}, & \Atomic(\text{Chiều}) &= \{3,4\}.
\end{align*}

Two timeslots overlap iff they occur on the same day and share at least one atomic block:
\begin{equation}
\Overlap(t_1, t_2) \iff day(t_1) = day(t_2) \wedge \bigl(\Atomic(period(t_1)) \cap \Atomic(period(t_2)) \neq \varnothing\bigr).
\label{eq:overlap}
\end{equation}

\section{Problem Definition}

The current system solves a \textbf{fixed-workbook personal timetable selection problem}. The workbook timetable is treated as immutable. The solver does not reschedule the institution-wide timetable. Instead, it selects a compatible personal timetable from existing LHP candidates.

Given a set of user-selected teaching-team requests $(course\_code, teaching\_team\_key)$, the system must:
\begin{enumerate}
    \item identify all valid teaching-unit candidates for each selected request,
    \item choose exactly one candidate teaching unit per selected request,
    \item ensure that the chosen LHPs do not overlap in time,
    \item enumerate up to $k_{\max}$ distinct feasible solutions, and
    \item rank feasible solutions by movement cost.
\end{enumerate}

This formulation is intentionally narrower than a global timetabling problem and matches the current product semantics.

\section{Decision Variables}

For every request $p \in \mathcal{P}$ and every candidate $c \in \mathcal{C}_p$, we introduce a Boolean variable:
\begin{equation}
x_{p,c} = \text{True} \iff \text{candidate } c \text{ is selected for request } p.
\label{eq:decision_var}
\end{equation}

The SAT model therefore represents a choice of one teaching-unit candidate per selected request.

\section{Core SAT Constraints}

\subsection{Exactly-One Candidate per Selected Request}

For each selected request, at least one candidate must be chosen:
\begin{equation}
\forall p \in \mathcal{P}: \bigvee_{c \in \mathcal{C}_p} x_{p,c}.
\label{eq:atleastone}
\end{equation}

For each selected request, at most one candidate may be chosen:
\begin{equation}
\forall p \in \mathcal{P},\; \forall c_1 < c_2 \in \mathcal{C}_p:
(\neg x_{p,c_1} \vee \neg x_{p,c_2}).
\label{eq:atmostone}
\end{equation}

Together, \eqref{eq:atleastone} and \eqref{eq:atmostone} enforce the exactly-one requirement.

\subsection{Cross-Request Conflict Constraints}

Two candidates from different selected requests are incompatible if either they reuse at least one session or any pair of their sessions overlaps in time.

Formally, for candidates $c \in \mathcal{C}_p$ and $d \in \mathcal{C}_q$ with $p \neq q$:
\begin{align}
\Conf(c,d) \iff &\; \bigl(\Sigma_{p,c} \cap \Sigma_{q,d} \neq \varnothing\bigr) \\
&\; \vee\; \exists s \in \Sigma_{p,c},\; s' \in \Sigma_{q,d}: \Overlap(t(s), t(s')).
\label{eq:conflict_def}
\end{align}

Whenever $\Conf(c,d)$ holds, the two candidates cannot be simultaneously selected:
\begin{equation}
\forall p \neq q,\; \forall c \in \mathcal{C}_p,\; \forall d \in \mathcal{C}_q:
\Conf(c,d) \Rightarrow (\neg x_{p,c} \vee \neg x_{q,d}).
\label{eq:conflict_clause}
\end{equation}

These clauses are sufficient to eliminate all pairwise timetable conflicts between chosen LHPs.

\section{Semantics of a SAT Model}

A satisfying assignment is interpreted as a personal timetable candidate if and only if:
\begin{itemize}
    \item for every selected request $p$, exactly one variable $x_{p,c}$ is assigned \texttt{True}, and
    \item no two chosen candidates violate the incompatibility clauses in \eqref{eq:conflict_clause}.
\end{itemize}

Thus, each SAT model corresponds to one feasible combination of LHP selections.

\section{Top-$k$ Enumeration by Blocking Clauses}

The solver enumerates multiple distinct solutions by repeatedly solving and then blocking the previous personal-timetable selection.

For a satisfying assignment, let $M^+ = \{v_1, \dots, v_r\}$ be precisely the
positive decision literals, with one selected candidate for every request in
$\mathcal{P}$. To exclude this exact combination in subsequent iterations, the
worker adds the blocking clause:
\begin{equation}
(\neg v_1 \vee \neg v_2 \vee \dots \vee \neg v_r).
\label{eq:blocking_clause}
\end{equation}

By the exactly-one constraints, this is equivalent to blocking the complete
personal-selection model while having length $|\mathcal{P}|$ rather than $N$.
It also remains valid if a future encoding introduces auxiliary SAT variables:
different values of auxiliary variables do not cause the same personal timetable
to be returned twice.

The procedure continues until one of the following conditions holds:
\begin{itemize}
    \item $k_{\max}$ feasible solutions have been collected,
    \item the remaining formula becomes UNSAT, or
    \item the global timeout is reached.
\end{itemize}

\section{Independent Personal-Timetable Validation}

The SAT encoding is followed by an independent validator. This validator defines the hard correctness conditions for the personal timetable output.

\subsection{Personal-Timetable Constraints (PT-1 to PT-7)}
\begin{itemize}
    \item \textbf{PT-1: Exactly one teaching unit per selected request.} Each requested $(course\_code, teaching\_team\_key)$ must produce exactly one chosen teaching unit.
    \item \textbf{PT-2: Candidate membership.} The chosen teaching unit must belong to the candidate set of the requested teaching-team request.
    \item \textbf{PT-3: Session-sequence consistency.} The returned session identifiers must exactly match the chosen candidate's canonical sequence, ordered by source row and then ordinal session identifier.
    \item \textbf{PT-4: Timeslot consistency.} The reported session timeslots must match the resolved timeslots in the fixed workbook schedule.
    \item \textbf{PT-5: No session reuse.} A session identifier cannot be reused across multiple selected teaching-team requests.
    \item \textbf{PT-6: No overlap among selected LHPs.} The chosen personal timetable must be pairwise non-overlapping under \eqref{eq:overlap}.
    \item \textbf{PT-7: No unexpected requests.} The output cannot contain a teaching-team request that was not explicitly selected by the user.
\end{itemize}

A personal timetable solution is considered valid only if all PT constraints hold.

\section{Workbook-Integrity Constraints}

In addition to personal-solution validation, the desktop application maintains a separate integrity layer for validating workbook-derived schedules. These constraints are denoted HC-1 to HC-6:
\begin{itemize}
    \item \textbf{HC-1: No physical room double-booking.}
    \item \textbf{HC-2: No overlap for individual lecturers.}
    \item \textbf{HC-3: No overlap for student cohorts.}
    \item \textbf{HC-4: Respect external lecturer blocks.}
    \item \textbf{HC-5: No internal LHP overlap for shared cohorts.}
    \item \textbf{HC-6: Lab-compatible room assignment for lab-like sessions.}
\end{itemize}

Conceptually, the HC layer validates timetable integrity, whereas the PT layer validates the correctness of the personal selector output.

\section{Soft Objective: Movement Cost Ranking}

Once hard-feasible personal timetables have been found, solutions are ranked by movement cost.

Consider the selected sessions on a fixed day, ordered chronologically. A transition contributes to the score only when two sessions are adjacent in atomic teaching blocks:
\begin{equation}
end(s_i) + 1 = start(s_{i+1}).
\label{eq:adjacency}
\end{equation}

For two adjacent sessions in rooms $r_i$ and $r_{i+1}$, the transition cost is defined as:
\begin{equation}
\Cost(r_i, r_{i+1}) =
\begin{cases}
0, & \text{same room}, \\
1, & \text{different rooms in the same building}, \\
2, & \text{different buildings in the same movement zone}, \\
3, & \text{different movement zones}.
\end{cases}
\label{eq:movement_cost}
\end{equation}

Special cases:
\begin{itemize}
    \item If either session is online, the transition cost is $0$.
    \item If either session has no physical room, the transition cost is $0$.
    \item If two sessions are not adjacent under \eqref{eq:adjacency}, no movement cost is added.
\end{itemize}

The total cost of a solution is therefore:
\begin{equation}
\text{MovementCost} = \sum_{\text{valid adjacent transitions}} \Cost(r_i, r_{i+1}).
\end{equation}

Solutions with lower movement cost are ranked ahead of solutions with higher cost.
The cost is evaluated after SAT feasibility and ranks only the bounded set of
solutions returned by enumeration. It is not optimized inside the CNF, so the
first returned solution is not necessarily the global movement optimum.

The room-building map used by the desktop implementation is:
\begin{itemize}
    \item buildings $A$ and $B$ map to zone $GD4$;
    \item buildings $T$ and \text{GĐ3} map to zone $GD3$;
    \item buildings $G2$, $G3$, $E3$, and $E5$ map to zone $GD2$;
    \item rooms containing \text{ĐHKHTN} map to the separate \text{ĐHKHTN} zone;
    \item an unmapped building name forms a zone with the same name.
\end{itemize}

\section{Encoding Complexity}

Let $m_p = |\mathcal{C}_p|$ and let
\begin{equation}
N = \sum_{p \in \mathcal{P}} m_p
\end{equation}
be the total number of candidate teaching units across all selected requests.

Then:
\begin{itemize}
    \item the at-least-one constraints contribute $|\mathcal{P}|$ clauses,
    \item the at-most-one constraints contribute $\sum_p \binom{m_p}{2}$ clauses,
    \item the cross-request conflict layer can contain $O(N^2)$ clauses in the worst case,
    \item conflict discovery is implemented with indexes from session identifiers and
    $(day, atomic\ period)$ buckets to candidate variables; its work is proportional
    to the occupied buckets rather than an unconditional all-pairs scan,
    \item each additional top-$k$ solution adds one blocking clause of the form \eqref{eq:blocking_clause}.
\end{itemize}

In practice, this formulation is substantially lighter than global timetable rescheduling because it only selects among fixed workbook candidates rather than searching over room-timeslot assignments for all sessions.

\section{Implementation Traceability}

The Windows desktop implementation corresponding to this specification is split across the following components:
\begin{itemize}
    \item candidate construction, SAT encoding, and movement-cost ranking:
    \path{scheduler/desktop/src/Scheduler.Application/}
    \item protocol v2 top-$k$ enumeration with exactly-one blocking groups:
    \path{scheduler/desktop/native/SolverWorker/}
    \item independent validation of HC and PT constraints:
    \path{scheduler/desktop/src/Scheduler.Application/TimetableValidator.cs}
\end{itemize}

The C\# selector and validator under \texttt{scheduler/desktop} are the product
implementation. Compatibility fixtures and the independent validator remain
release evidence for the desktop implementation.

\end{document}
